\section*{Hoofdstuk \ref{ch:b2:equilibrium-shifts} --- Evenwichtsconstanten en evenwichtsverschuivingen}

\begin{solution}{exo:b2:equilibrium-shifts:1}
$K^\circ = \exp(32\,800/(8.314 \times 298.15)) = \num{5.6e5}$, en voor
$+\qty{32.8}{kJ/mol}$ het omgekeerde, \num{1.8e-6}. Een factor 100 vereist
$\Delta(\Delta_r G^\circ) = -RT\ln 100 = \qty{-11.4}{kJ/mol}$.
\end{solution}

\begin{solution}{exo:b2:equilibrium-shifts:2}
$\ln K^\circ(400) = \ln(5.4 \times 10^5) + (91\,880/8.314)(1/400 -
1/298.15) = 13.20 - 9.44 = 3.76$, $K^\circ \approx 43$. De tabellen geven 35.6:
over \qty{100}{K} zit de benadering met constante $\Delta_r H^\circ$ er al
20\,\% naast (de reactie wordt exothermer naarmate $T$ stijgt).
\end{solution}

\begin{solution}{exo:b2:equilibrium-shifts:3}
(a) $n = 1$, $r = 0$, $\varphi = 2$: $v = 1$ (de dampdruk is een functie
van $T$). (b) $v = 3 - 1 + 2 - 3 = 1$. (c) $v = 3 - 1 + 2 - 1 = 3$. (d) Vier
deeltjes, één reactie, twee fasen (de koolstof en het gas): $v = 4 - 1 + 2 - 2 =
3$; als het gas alleen uit koolstof en stoom gemaakt wordt, is $n(\ce{CO}) = n(\ce{H2})$
in het gas, $s = 1$ en $v = 2$.
\end{solution}

\begin{solution}{exo:b2:equilibrium-shifts:4}
Verwarmen: \omterm{def:b2:reaction-enthalpy:exothermic}{endotherme} richting, naar \ce{NO2} (voorwaarts). Samenpersen:
naar minder gasmoleculen, \ce{N2O4} (achterwaarts). Argon bij constant volume:
geen verandering, de partiële drukken van de reagerende gassen blijven gelijk.
\end{solution}

\begin{solution}{exo:b2:equilibrium-shifts:5}
Bij \qty{298}{K}: $\Delta_r G^\circ = 180.6 - 298.15 \times 0.0248 =
\qty{173.2}{kJ/mol}$, $K^\circ = \num{4.5e-31}$. Bij \qty{2000}{K}: $\Delta_r
G^\circ = 180.6 - 2000 \times 0.0248 = \qty{131.0}{kJ/mol}$, $K^\circ =
\num{3.8e-4}$. Met $\Delta\nu_{\text{gas}} = 0$ is $x(\ce{NO})^2 = K^\circ x(\ce{N2})
x(\ce{O2})$: $x(\ce{NO}) = \sqrt{3.8 \times 10^{-4} \times 0.78 \times 0.21} =
\num{7.9e-3}$, bijna één procent. Motoren verbranden bij zulke temperaturen; terwijl de
gassen afkoelen, daalt $K^\circ$ enorm, maar de ontleding van \ce{NO} is
zeer traag: het hete evenwicht wordt in de uitlaatgassen bevroren (tot een katalysator
het verwijdert).
\end{solution}

\begin{solution}{exo:b2:equilibrium-shifts:6}
Met $\xi$ mol ester en een constante totale hoeveelheid is $Q =
\xi^2/[(1 - \xi)(n_B - \xi)]$. Voor $n_B = 1.0$: $\xi^2/(1 - \xi)^2 = 4$, $\xi
= \qty{0.67}{mol}$. Voor $n_B = 3.0$: $3\xi^2 - 16\xi + 12 = 0$, $\xi =
\qty{0.90}{mol}$. Water verwijderen houdt $Q$ onder $K^\circ$, wat $\xi$ ook is: de
reactie blijft voorwaarts verlopen tot het zuur verbruikt is.
\end{solution}

\begin{solution}{exo:b2:equilibrium-shifts:7}
Zonder argon: hoeveelheden $1 - \alpha$, $2\alpha$, totaal $1 + \alpha$;
$Q = 4\alpha^2/(1 - \alpha^2) = 0.148$, $\alpha = \sqrt{0.148/4.148} = 0.189$.
Met \qty{1}{mol} argon bij \qty{1}{bar}: totaal $2 + \alpha$, $Q =
4\alpha^2/[(1 - \alpha)(2 + \alpha)] = 0.148$, dus $4.148\alpha^2 + 0.148\alpha
- 0.296 = 0$, $\alpha = 0.250$: verdunnen bij constante totale druk verlaagt de
partiële drukken, wat de kant met meer gasmoleculen begunstigt. Bij
constant volume veranderen de partiële drukken van de reagerende gassen niet:
$\alpha$ blijft 0.189.
\end{solution}

\begin{solution}{exo:b2:equilibrium-shifts:8}
Alleen de vaste stof: $n = 3$, $r = 1$, $s = 1$ ($p(\ce{NH3}) = p(\ce{HCl})$),
$\varphi = 2$: $v = 3 - 1 - 1 + 2 - 2 = 1$; kies $T$ en de druk ligt
vast (een ``dissociatiedruk''). Met toegevoegde ammoniak is $s = 0$ en $v = 2$:
$T$ en één druk kunnen gekozen worden; het product $p(\ce{NH3})\,p(\ce{HCl})$ ligt
vast door $T$.
\end{solution}

\begin{solution}{exo:b2:equilibrium-shifts:9}
$\Delta_r H^\circ = -R\ln(K_2/K_1)/(1/T_2 - 1/T_1) = -8.314 \times
\ln(0.0173)/(1/600 - 1/500) = \qty{-101}{kJ/mol}$, negatiever dan bij
\qty{298}{K} (\qty{-91.9}{kJ/mol}), zoals de \omterm{thm:b2:reaction-enthalpy:kirchhoff}{wet van Kirchhoff} voorspelde
($\Delta_r C_p^\circ < 0$).
\end{solution}

\begin{solution}{exo:b2:equilibrium-shifts:10}
$Q = \dfrac{n_{\ce{NH3}}^2\,n_{\text{tot}}^2}{n_{\ce{N2}}n_{\ce{H2}}^3}
\left(\dfrac{p^\circ}{p}\right)^2$. Bij constante $T$, $p$:
$\partial\ln Q/\partial n_{\ce{N2}} = -1/n_{\ce{N2}} + 2/n_{\text{tot}}$, positief
wanneer $n_{\ce{N2}} > n_{\text{tot}}/2$, dus $x(\ce{N2}) > 1/2$. Dan wordt $Q$
groter dan $K^\circ$ en verschuift het evenwicht achterwaarts. Bij constant volume is
$\partial\ln Q/\partial n_{\ce{N2}} = -1/n_{\ce{N2}} < 0$: altijd voorwaarts.
\end{solution}

\begin{solution}{exo:b2:equilibrium-shifts:11}
$K^\circ = \exp(-1620/(8.314 \times 1100)) = 0.84$: $p(\ce{CO2}) =
\qty{0.84}{bar}$ bij evenwicht, dus $n = pV/RT = 0.84 \times 10^5 \times
0.0100/(8.314 \times 1100) = \qty{0.092}{mol}$ gas. Met \qty{0.050}{mol}
kalksteen ontleedt alles ($p = \qty{0.46}{bar} < K^\circ p^\circ$):
geen evenwicht, het evenwicht wordt verbroken. Met \qty{0.200}{mol} evenwicht:
\qty{0.092}{mol} \ce{CO2} en \ce{CaO}, \qty{0.108}{mol} \ce{CaCO3} over.
\end{solution}

\begin{solution}{exo:b2:equilibrium-shifts:12}
Bed 1 eindigt bij ongeveer \qty{890}{K} en 68\,\% omzetting, bed 2 bij ongeveer
\qty{785}{K} en 92\,\%, bed 3 bij ongeveer \qty{725}{K} en 98\,\%. In één
adiabatisch bed warmt het gas op terwijl het reageert, tot het de evenwichtskromme
bereikt, die bij die temperatuur slechts ongeveer 68\,\% toelaat. Een overal koude katalysator
is geen optie, omdat de katalysator alleen heet werkt; afkoelen
tussen de bedden behoudt de snelheid en verlegt de evenwichtsgrens omhoog. Het laatste bed
wordt begrensd door het evenwicht bij de laagste temperatuur waarbij de katalysator
nog werkt (fabrieken absorberen dan het trioxide en leiden het gas nog eens over
een katalysator).
\end{solution}

\begin{solution}{pb:b2:equilibrium-shifts:1}
\textbf{1.} $\Delta_r H^\circ = \qty{-91.88}{kJ/mol}$, $\Delta_r S^\circ =
2(192.77) - 191.61 - 3(130.68) = \qty{-198.1}{J/(K.mol)}$, $\Delta_r G^\circ =
\qty{-32.8}{kJ/mol}$, $K^\circ = \num{5.6e5}$.
\textbf{2.} $K^\circ(700) = \exp(-54\,380/(8.314 \times 700)) = \num{8.8e-5}$;
$K^\circ(800) = \exp(-77\,320/(8.314 \times 800)) = \num{8.9e-6}$.
\textbf{3.} $\Delta_r H^\circ = -R\ln(K_{800}/K_{700})/(1/800 - 1/700)$,
dus \qty{-106}{kJ/mol}, negatiever dan bij \qty{298}{K}.
\textbf{4.} $\ln K^\circ(723) = \ln(8.75 \times 10^{-5}) + 12\,777 \times
(1/723.15 - 1/700)$, met $12\,777 = 106\,230/8.314$: $K^\circ =
\num{4.9e-5}$.
\textbf{5.} $\Delta_r G^\circ(723) = -91.88 + 723.15 \times 0.1981 =
\qty{51.4}{kJ/mol}$, $K^\circ = \num{1.9e-4}$: viermaal te groot, omdat
$\Delta_r H^\circ$ en $\Delta_r S^\circ$ over \qty{425}{K} veranderen
($\Delta_r C_p^\circ \approx \qty{-44}{J/(K.mol)}$).
\textbf{6.} \ce{N2} $1 - \xi$, \ce{H2} $3 - 3\xi$, \ce{NH3} $2\xi$, totaal $4 -
2\xi$.
\textbf{7.} $x = 2\xi/(4 - 2\xi)$, dus $1 - x = (4 - 4\xi)/(4 - 2\xi)$; de
stikstoffractie is $(1 - \xi)/(4 - 2\xi) = (1 - x)/4$ en de waterstoffractie
is driemaal groter.
\textbf{8.}
\[
  K^\circ = \frac{x^2}{\frac{1-x}{4}\left(\frac{3(1-x)}{4}\right)^3}
  \left(\frac{p^\circ}{p}\right)^2 = \frac{256}{27}\,\frac{x^2}{(1 - x)^4}
  \left(\frac{p^\circ}{p}\right)^2 ;
\]
neem de vierkantswortel.
\textbf{9.} $a = (\sqrt{27}/16) \times \sqrt{4.9 \times 10^{-5}} \times 1 =
\num{2.27e-3}$; $x/(1 - x)^2 = a$ geeft $x = 0.0023$.
\textbf{10.} $a = 0.455$; $ax^2 - (2a + 1)x + a = 0$: $x = 0.25$.
\textbf{11.} $n = 3$, $r = 1$, $s = 1$ (verhouding 1 : 3, door de reactie behouden),
$\varphi = 1$: $v = 2$. Zodra $T$ en $p$ gekozen zijn, blijft er niets vrij.
\textbf{12.} Druk: voorwaarts ($\Delta\nu_{\text{gas}} = -2$). Temperatuur:
achterwaarts (\omterm{def:b2:reaction-enthalpy:exothermic}{exotherm}).
\textbf{13.} Een inert gas bij constante totale druk verlaagt de partiële
drukken van de reagerende gassen, zoals een drukdaling zou doen: achterwaarts.
\textbf{14.} $\partial\ln Q/\partial n_{\ce{N2}} = (-1/0.30 + 2)/n_{\text{tot}} < 0$:
$Q$ zakt onder $K^\circ$, voorwaarts.
\textbf{15.} Nee: bij constant volume is $\partial\ln Q/\partial n_{\ce{N2}} =
-1/n_{\ce{N2}} < 0$, wat de samenstelling ook is: voorwaarts.
\textbf{16.} Met bijna geen ammoniak ligt $Q$ aan de inlaat ver onder $K^\circ$:
het gas komt ver van het evenwicht binnen en reageert voorwaarts.
\textbf{17.} Het gas verblijft te kort op de katalysator om het
evenwicht te bereiken; een langer contact zou een grotere, duurdere
converter kosten voor een kleine winst.
\textbf{18.} $x = 2\xi/(4 - 2\xi) = 0.18$ geeft $\xi = \qty{0.305}{mol}$ per
mol gevoede \ce{N2}: \omterm{def:b2:equilibrium-shifts:single-pass}{conversie per doorgang} van waterstof (en van stikstof)
$0.305$, 31\,\%.
\textbf{19.} In stationaire werking wordt alle verse voeding uiteindelijk omgezet:
per tijdseenheid gaat er \qty{2.00}{mol} ammoniak naar buiten; de converter moet per tijdseenheid $1.00/
0.305 = \qty{3.3}{mol}$ \ce{N2} doorlaten, de rest wordt teruggevoerd.
\textbf{20.} Het argon (en het methaan) dat met de voeding binnenkomt, reageert nooit en
wordt niet met de vloeibare ammoniak afgevoerd: zonder \omterm{def:b2:equilibrium-shifts:single-pass}{spuistroom} zou het zich
ophopen in de kringloop en de partiële drukken van de beginstoffen verlagen.
\textbf{21.} Bij \qty{500}{K} is de ijzerkatalysator te traag: het evenwicht
zou gunstig zijn, maar in een industriële tijdspanne nooit benaderd worden.
\textbf{22.} Bij \qty{723}{K} en \qty{200}{bar}, stoichiometrische voeding, ideale
gassen: $\boldsymbol{x(\ce{NH3}) \approx 0.25}$.
\end{solution}
